\documentclass[../../../include/open-logic-section]{subfiles}

\begin{document}

\olfileid[id]{sth}{ordinals}{basic}
\olsection{Sifat-Sifat Dasar Ordinal}

Kita telah mengamati bahwa beberapa ordinal pertama merupakan bilangan asli.
Alasan utama untuk mengembangkan teori ordinal adalah memperluas prinsip
induksi yang berlaku pada bilangan asli. Kita akan menuju hasil ini melalui
serangkaian hasil elementer.

\begin{lem}\ollabel{ordmemberord}
Setiap !!{element} dari suatu ordinal merupakan suatu ordinal.
\end{lem}

\begin{proof}
Misalkan $\alpha$ suatu ordinal dengan $b \in \alpha$. Karena $\alpha$
bersifat transitif, $b \subseteq \alpha$. Jadi, $\in$ mengurutkan $b$ dengan
baik karena $\in$ mengurutkan $\alpha$ dengan baik.

Untuk melihat bahwa $b$ bersifat transitif, andaikan $x \in c \in b$. Jadi,
$c \in \alpha$ karena $b \subseteq \alpha$. Sekali lagi, karena $\alpha$
bersifat transitif, $c \subseteq \alpha$, sehingga $x \in \alpha$. Jadi $x,
c, b \in \alpha$. Namun, $\in$ mengurutkan $\alpha$ dengan baik, sehingga
$\in$ merupakan relasi transitif pada $\alpha$ menurut
\olref[sth][ordinals][wo]{wo:strictorder}. Karena itu, dari $x \in c \in b$
kita memperoleh $x \in b$. Dengan menggeneralisasi, $c \subseteq b$.
\end{proof}

\begin{cor}\ollabel{ordissetofsmallerord}
$\alpha = \Setabs{\beta \in \alpha}{\beta \text{ merupakan suatu ordinal}}$,
untuk setiap ordinal~$\alpha$.
\end{cor}

\begin{proof}
Langsung dari \olref{ordmemberord}.
\end{proof}

Gagasan kasar di balik dua hasil utama berikut,
\olref{ordinductionschema} dan \olref{ordtrichotomy}, adalah bahwa
ordinal-ordinal itu sendiri diurutkan dengan baik oleh relasi keanggotaan:

\begin{thm}[Induksi Transfinit]\ollabel{ordinductionschema}
Untuk sembarang formula $\phi(x)$:
\[
	\text{jika }\exists \alpha \phi(\alpha)\text{, maka }\exists \alpha(\phi(\alpha)
	\land  (\forall \beta \in \alpha) \lnot \phi(\beta))
\]
dengan kuantor-kuantor yang ditampilkan secara implisit dibatasi pada
ordinal.
\end{thm}

\begin{proof}
%\olref{ordinductionschema1}.
Andaikan $\phi(\alpha)$ untuk suatu ordinal $\alpha$. Jika $(\forall \beta
\in \alpha) \lnot \phi(\beta)$, maka pembuktian selesai. Jika tidak, karena
$\alpha$ merupakan suatu ordinal, terdapat suatu !!{element} terkecil
menurut $\in$ yang memenuhi $\phi$, dan anggota ini merupakan
suatu ordinal menurut \olref{ordmemberord}.
%	\olref{ordinductionschema2}. Suppose there is some ordinal
%	$\gamma$ such that $\lnot\phi(\gamma)$. Then by
%	\olref{ordinductionschema1}, there is an $\in$-minimal ordinal
%	$\alpha$ for which $\lnot\phi(\alpha)$. So $(\forall \beta <
%	\alpha) \phi(\beta)$, rendering the antecedent of the conditional
%	false.
\end{proof}
\noindent
Perhatikan bahwa kita juga dapat mengungkapkan \olref{ordinductionschema}
sebagai skema:
\[
\text{jika }\forall \alpha((\forall \beta \in \alpha)\phi(\beta) \lif
\phi(\alpha))\text{, maka }\forall \alpha\phi(\alpha)
\]
cukup dengan mengambil $\lnot\phi(\alpha)$ dalam
\olref{ordinductionschema}, lalu melakukan manipulasi logis elementer.

\begin{thm}[Trikotomi]\ollabel{ordtrichotomy}
$\alpha \in \beta \lor \alpha = \beta \lor \beta \in \alpha$, untuk
sembarang ordinal $\alpha$ dan $\beta$.
\end{thm}

\begin{proof}
Pembuktian dilakukan dengan induksi ganda, yakni dengan menggunakan
\olref{ordinductionschema} dua kali. Katakan bahwa $x$ \emph{dapat
dibandingkan} dengan $y$ jika dan hanya jika $x \in y \lor x = y \lor y \in
x$.

Untuk induksi, andaikan setiap ordinal dalam~$\alpha$ dapat dibandingkan
dengan \emph{setiap} ordinal. Untuk induksi berikutnya, andaikan $\alpha$
dapat dibandingkan dengan setiap ordinal dalam~$\beta$. Kita akan
menunjukkan bahwa $\alpha$ dapat dibandingkan dengan~$\beta$. Dengan induksi
pada~$\beta$, akan mengikuti bahwa $\alpha$ dapat dibandingkan dengan setiap
ordinal; kemudian, dengan induksi pada~$\alpha$, \emph{setiap} ordinal dapat
dibandingkan dengan \emph{setiap} ordinal, sebagaimana diperlukan. Cukup
andaikan bahwa $\alpha \notin \beta$ dan $\beta \notin \alpha$, lalu
tunjukkan bahwa $\alpha = \beta$.

Untuk menunjukkan bahwa $\alpha \subseteq \beta$, tetapkan $\gamma \in
\alpha$; ini merupakan suatu ordinal menurut \olref{ordmemberord}. Jadi,
menurut hipotesis induksi pertama, $\gamma$ dapat dibandingkan dengan
$\beta$. Namun, jika $\gamma = \beta$ ataupun $\beta \in \gamma$, maka
$\beta \in \alpha$ (dengan menggunakan fakta bahwa $\alpha$ bersifat
transitif jika diperlukan), berlawanan dengan pengandaian kita; jadi $\gamma
\in \beta$. Dengan menggeneralisasi, $\alpha \subseteq \beta$.

Penalaran yang persis serupa, dengan menggunakan hipotesis induksi kedua,
menunjukkan bahwa $\beta \subseteq \alpha$. Jadi $\alpha = \beta$.
\end{proof}\noindent
Dengan demikian, terkadang kita akan menulis $\alpha <\beta$ alih-alih
$\alpha \in \beta$, karena $\in$ berperilaku sebagai relasi pengurutan. Tidak
ada alasan mendalam untuk hal ini, selain kelazimannya dan karena lebih mudah
menulis $\alpha \leq \beta$ daripada $\alpha \in \beta \lor \alpha =
\beta$.\footnote{Kita dapat menulis $\alpha
\mathrel{\underline{\in}} \beta$; tetapi notasi itu sama sekali tidak
standar.}

Berikut dua konsekuensi singkat dari hasil-hasil terakhir kita; yang pertama
menerapkan notasi baru kita:

\begin{cor}\ollabel{ordordered}
Jika $\exists \alpha\phi(\alpha)$, maka $\exists \alpha(\phi(\alpha)
\land \forall \beta(\phi(\beta) \lif \alpha \leq \beta))$.
Selain itu, untuk sembarang ordinal $\alpha, \beta, \gamma$, berlaku baik
$\alpha \notin \alpha$ maupun $\alpha \in \beta \in \gamma \lif \alpha
\in \gamma$.
\end{cor}

\begin{proof}
Sama seperti \olref[wo]{wo:strictorder}.
\end{proof}

\begin{prob}
Lengkapi ``penalaran yang persis serupa'' dalam pembuktian
\olref[sth][ordinals][basic]{ordtrichotomy}.
\end{prob}

\begin{cor}\ollabel{corordtransitiveord}
$A$ merupakan suatu ordinal jika dan hanya jika $A$ merupakan suatu himpunan
transitif dari ordinal-ordinal.
\end{cor}

\begin{proof}
\emph{Dari kiri ke kanan.} Menurut \olref{ordmemberord}.
\emph{Dari kanan ke kiri.} Jika $A$ merupakan suatu himpunan transitif dari
ordinal-ordinal, maka $\in$ mengurutkan $A$ dengan baik menurut
\olref{ordinductionschema} dan \olref{ordtrichotomy}.
\end{proof}

Di atas, kita menyarikan
\olref{ordinductionschema} dan \olref{ordtrichotomy} sebagai pernyataan bahwa
$\in$ mengurutkan ordinal-ordinal dengan baik. Namun, kita harus
\emph{sangat berhati-hati} terhadap pernyataan semacam ini karena hasil
berikut:

\begin{thm}[Paradoks Burali--Forti]\ollabel{buraliforti}
Tidak ada himpunan semua ordinal.
\end{thm}

\begin{proof}
Untuk reductio, andaikan $O$ merupakan himpunan semua ordinal. Jika $\alpha
\in \beta \in O$, maka $\alpha$ merupakan suatu ordinal menurut
\olref{ordmemberord}, sehingga $\alpha \in O$. Jadi $O$ bersifat transitif,
dan karena itu $O$ merupakan suatu ordinal menurut
\olref{corordtransitiveord}. Maka $O \in O$, bertentangan dengan
\olref{ordordered}.
\end{proof}
\noindent
Hasil ini dinamai menurut \citeauthor{Burali-Forti1897}. Namun, Cantor-lah
yang pada tahun 1899---dalam sebuah surat kepada Dedekind---pertama kali
melihat dengan jelas \emph{kontradiksi} dalam pengandaian bahwa terdapat
suatu himpunan semua ordinal. Sebagaimana dijelaskan van Heijenoort:
\begin{quote}
  Burali-Forti sendiri menganggap kontradiksi tersebut menetapkan, melalui
  \emph{reductio ad absurdum}, hasil bahwa pengurutan alami ordinal hanyalah
  suatu urutan parsial.
  \citep[p.~105]{Heijenoort1967}
\end{quote}
Dengan mengesampingkan kekeliruan Burali-Forti, kita dapat merangkum uraian
di atas sebagai berikut. Ordinal merupakan himpunan-himpunan yang secara
individual diurutkan dengan baik oleh keanggotaan, dan secara kolektif
diurutkan dengan baik oleh keanggotaan (tanpa secara kolektif membentuk suatu
himpunan).

Sebagai penutup, berikut beberapa sifat dasar lain tentang ordinal yang
mengikuti dari \olref{ordinductionschema} dan \olref{ordtrichotomy}.

\begin{prop}
Setiap barisan ordinal yang menurun secara ketat bersifat hingga.
\end{prop}

\begin{proof}
Setiap barisan ordinal tak hingga yang menurun secara ketat $\alpha_0 >
\alpha_1 > \alpha_2 > \ldots$ tidak mempunyai anggota minimal menurut $<$,
bertentangan dengan \olref{ordinductionschema}.
\end{proof}

\begin{prop}\ollabel{ordinalsaresubsets}
$\alpha \subseteq \beta \lor \beta \subseteq \alpha$, untuk sembarang
ordinal $\alpha, \beta$.
\end{prop}

\begin{proof}
Jika $\alpha \in \beta$, maka $\alpha \subseteq \beta$ karena $\beta$
bersifat transitif. Serupa dengan itu, jika $\beta \in \alpha$, maka $\beta
\subseteq \alpha$. Jika $\alpha = \beta$, maka $\alpha \subseteq \beta$ dan
$\beta \subseteq \alpha$. Jadi, menurut \olref{ordtrichotomy}, pembuktian
selesai.
\end{proof}

\begin{prop}\ollabel{ordisoidentity}
$\alpha = \beta$ jika dan hanya jika $\ordeq{\alpha}{\beta}$, untuk
sembarang ordinal $\alpha, \beta$.
\end{prop}

\begin{proof}
Ordinal-ordinal merupakan pengurutan baik; jadi, hasil ini langsung mengikuti
dari Trikotomi (\olref[basic]{ordtrichotomy}) dan
\olref[iso]{wellordnotinitial}.
\end{proof}

\begin{prob}\ollabel{probunionordinalsordinal}
Buktikan bahwa, jika setiap anggota $X$ merupakan suatu ordinal, maka
$\bigcup X$ merupakan suatu ordinal.
\end{prob}

\end{document}
